\begin{proposition}[Exclusion of Inoue general fibres]\hypertarget{tag-01F7}{}\label{cdp:st:C.3.12}
\StatementRecord{cdp:st:C.3.12}{proved}{01F7}{\SeriesRef{cdp:st:C.3.1}, \SeriesRef{cdp:eq:C.IN.2}, \SeriesRef{cdp:eq:C.IN.3}, \SeriesRef{cdp:eq:C.IN.6}; \cite[pp.~274, 278, 280]{Inoue1974}; \cite[Sections 2.1--2.3]{KhaledTeleman2026}; \cite[Theorem 3.2]{BMS2012}; \cite[Section 3.3]{BMS2012}}


Under the hypotheses of \ref{cdp:st:C.3.1}, the very general smooth fibre is
not an Inoue surface. More precisely, a proper holomorphic map from
a compact connected complex manifold to a compact Riemann surface,
whose fibres over the complement of a finite set are Inoue surfaces,
has a class in $H^1(Y;\mathbb R)$ restricting nontrivially to each
such fibre. Thus these families cannot have $b_1(Y)=0$.
This assertion permits arbitrary multiplicities and arbitrary
singularities of the omitted fibres.
\end{proposition}


\begin{proof}

\begin{enumerate}[leftmargin=*]
\item Write $B^\circ$ for the complement of the
finite set and $Y^\circ=f^{-1}(B^\circ)$. Put
$\mathbb H=\{w\in\mathbb C:\operatorname{Im}w>0\}$.
Each Inoue surface is a quotient of $\mathbb H\times\mathbb C$.
In the three standard presentations the first coordinates of all
deck transformations have the form
\begin{equation}
 w\longmapsto a w+b,\qquad a>0,\quad b\in\mathbb R.
\hypertarget{tag-00NK}{}\label{cdp:eq:C.IN.1}
\end{equation}
They include a nonzero real translation and a transformation
$w\mapsto\alpha w$, with $\alpha>1$. Inoue's original
Propositions 2 and 5 establish infinitesimal rigidity for types
$S_M$ and $S^{-}$; his Proposition 4 gives a locally complete
one-parameter family for type $S^{+}$
\cite[pp.~274, 278, 280]{Inoue1974}.
In that family the parameter changes only the second coordinate:
local coordinate transitions have the form
\begin{equation}
\begin{aligned}
 w_i&=\alpha^{m_{ij}}w_j+a_{ij},\\
 z_i&=z_j+b_{ij}w_j+c_{ij}+t\,m_{ij},
\end{aligned}
\hypertarget{tag-00NL}{}\label{cdp:eq:C.IN.2}
\end{equation}
where $m_{ij}\in\mathbb Z$ and the other coefficients are
independent of $t$. In particular the real affine transformations
in the first coordinate are constant throughout this complete
family. The standard deck presentations are also recorded in
\cite[Sections 2.1--2.3]{KhaledTeleman2026}.

\item After shrinking a simply connected base disc $V\subset B^\circ$,
local completeness identifies $f^{-1}(V)$ with the pullback of one
of these families. Its universal cover over $V$ is therefore
$\mathbb H\times\mathbb C\times V$, with holomorphic covering
map to $f^{-1}(V)$. In the $S^{+}$ case the parameter $t$ in
\eqref{cdp:eq:C.IN.2} is a holomorphic function on $V$; the other two cases
are products. The smooth real one-form
\begin{equation}
 \theta_V=d\log\operatorname{Im}w
\hypertarget{tag-00NM}{}\label{cdp:eq:C.IN.3}
\end{equation}
is invariant under every deck transformation, where $d$ is the
exterior derivative on the \emph{whole} covering space. It thus
descends to a closed real one-form on $f^{-1}(V)$.

\item Over a simply connected smaller overlap, a transition
between the two family descriptions lifts holomorphically to the
universal covers. On each fibre its first coordinate is independent
of $z$: composing with a Cayley transform gives a bounded entire
function of $z$, which is constant. Applying the same argument to
the inverse proves that the resulting first coordinate is an
automorphism $h$ of $\mathbb H$. This automorphism conjugates the
first-coordinate deck groups. Their common fixed-point set on
$\partial\mathbb H=\mathbb R\cup\{\infty\}$ is exactly
$\{\infty\}$: the nonzero translation has this as its sole fixed
point, and all their elements fix it. Consequently $h(\infty)=
\infty$ and $h(w)=a w+b$ with $a>0$ and $b\in\mathbb R$.
In a holomorphic family of these transitions, $a$ and $b$ are
holomorphic functions of the base parameter; this follows, for
example, by taking the $w$ derivative and then evaluating at a
fixed $w$. Their real values force them to be locally constant.
Therefore
\begin{equation}
 d\log\operatorname{Im}(a w+b)
 =d\log(a\operatorname{Im}w)
 =d\log\operatorname{Im}w
\hypertarget{tag-00NN}{}\label{cdp:eq:C.IN.4}
\end{equation}
as forms on the total family. Thus the $\theta_V$ glue to an
actual closed form $\theta$ on $Y^\circ$.
The same argument includes a finite \'etale quotient of an Inoue
surface. Locally over a simply connected base disc its finite
cover extends to the family: lift a differentiable marking by
Ehresmann's theorem and equip the resulting covering space with
the pulled-back complex structure. Passing to the Galois closure
when necessary gives a family of finite covers of Inoue surfaces;
these are again Inoue surfaces, as is also seen from their
finite-index lattice presentations. Every deck transformation of
the quotient normalizes a finite-index Inoue deck group on
$\mathbb H\times\mathbb C$. Its first-coordinate group still
contains a nonzero translation and has common boundary fixed
point precisely $\infty$. The preceding conjugation argument
therefore makes each quotient deck transformation real affine in
$w$, and its coefficients are constant over the base disc.
Consequently \eqref{cdp:eq:C.IN.3} descends through the finite quotient as
well. Its pullback to the Inoue cover has the nonzero period
$\log\alpha$, so the descended fibre class is nonzero.
For a fibre loop represented on its universal cover by the deck
transformation $w\mapsto\alpha w$, its integral is
$\log\alpha\ne0$. Hence $[\theta]$ restricts nontrivially to
the fibre.

\item Put
$D=Y\setminus Y^\circ$. Let $\iota:\Delta\to Y$ be a small
holomorphic transverse disc at a smooth point of a reduced
irreducible component of $D$, avoiding all other components.
The pullback of the preceding transverse coordinate atlas on
$\Delta^*$ has constant real affine transitions. On the universal
cover
\[
 \mathbb H_c=\{\operatorname{Im}\tau>c\}
 \longrightarrow\Delta^*,\qquad
 \tau\longmapsto\exp(2\pi\mathrm i\tau),
\]
analytic continuation consequently gives a holomorphic developing
map $W:\mathbb H_c\to\mathbb H$ with
$W(\tau+1)=aW(\tau)+b$ for some $a>0$, $b\in\mathbb R$.
This follows by trivializing the pulled-back constant transition
cocycle on the simply connected covering space.
The period of $\iota^*\theta$ on the positive puncture loop is
$\log a$.

\item Use the Poincar\'e metrics
$|dw|/\operatorname{Im}w$ and
$|d\tau|/(\operatorname{Im}\tau-c)$.
The function $\log\operatorname{Im}w$ has differential of norm
one for the first metric. Schwarz--Pick and the horizontal segment
from $\tau$ to $\tau+1$ therefore give
\begin{equation}
\begin{aligned}
 |\log a|
 &\le d_{\mathbb H}\bigl(W(\tau),W(\tau+1)\bigr)\\
 &\le d_{\mathbb H_c}(\tau,\tau+1)
 \le\frac{1}{\operatorname{Im}\tau-c}.
\end{aligned}
\hypertarget{tag-00NO}{}\label{cdp:eq:C.IN.5}
\end{equation}
Letting $\operatorname{Im}\tau\to+\infty$ proves $a=1$.
Every meridian period of $\theta$ is thus zero. This calculation
applies to the reduced divisor and is unaffected by the
scheme-theoretic multiplicities of its components.

\item The map $H_1(Y^\circ;\mathbb R)\to H_1(Y;\mathbb R)$ is
surjective: loops can be moved away from a complex divisor.
Its kernel is generated by the meridians of the irreducible
components of $D$. Indeed, a bounding real two-chain can be put
in general position, transverse to the smooth divisor locus and
disjoint from its singular locus, which has real codimension at
least four. Removing small discs about its intersection points
leaves a chain in $Y^\circ$ whose additional boundary circles
are meridians. The period homomorphism of $\theta$ therefore factors uniquely
through $H_1(Y;\mathbb R)$ and defines a class
\begin{equation}
 \widetilde\theta\in H^1(Y;\mathbb R),\qquad
 \widetilde\theta|_{Y^\circ}=[\theta].
\hypertarget{tag-00NP}{}\label{cdp:eq:C.IN.6}
\end{equation}
Its restriction to an Inoue fibre is nonzero by \eqref{cdp:eq:C.IN.3}.
Consequently $b_1(Y)>0$.

\item Suppose a very general smooth fibre under the hypotheses of
\ref{cdp:st:C.3.1} is an Inoue surface. Ehresmann's
theorem identifies the diffeomorphism type, and in particular
the fundamental group, of all smooth fibres. The first two
Betti numbers are then $(1,0)$. These surfaces are minimal by
\ref{cdp:st:C.3.1} and non-K\"ahler. The compact-surface classification
gives class VII surfaces or non-K\"ahler elliptic surfaces
\cite[Theorem 3.2]{BMS2012}.
The Inoue group is solvable and is not virtually nilpotent.
For the latter assertion, its normal lattice is $\mathbb Z^3$
in type $S_M$, and a three-dimensional integral Heisenberg group
in types $S^{+}$ and $S^{-}$. Conjugation by a generator of the
cyclic quotient has an eigenvalue $\alpha>1$ on the rational
abelianization of this lattice. Every finite-index subgroup
contains a positive power of that generator and a finite-index
sublattice, so the eigenvalue becomes $\alpha^k>1$. In a
nilpotent group repeated commutators make every such conjugation
unipotent on this rational abelianization. 
\item A Hopf group is virtually cyclic, and a Kodaira group is virtually
nilpotent; neither can equal this group. For a properly elliptic
surface of Euler number zero the canonical bundle formula makes
the base orbifold hyperbolic: its orbifold canonical degree is
positive exactly in Kodaira dimension one. The elliptic
fibration induces a surjection of its fundamental group onto the
orbifold group. That group has a finite-index closed surface
group of genus at least two, and hence is not solvable. This
also cannot be a quotient of an Inoue group. The orbifold exact
sequence is the one used in \cite[Section 3.3]{BMS2012}.
Thus every smooth fibre is an Inoue surface. Applying \eqref{cdp:eq:C.IN.6}
contradicts $b_1(Y)=0$.

\end{enumerate}
\end{proof}